\documentclass[krantz1]{krantz}
\usepackage{fixltx2e,fix-cm}
\usepackage{amssymb, bbm, bm, mathtools}
\usepackage{amsmath}
\usepackage{graphicx}
\usepackage{color, rotating}
\usepackage{subcaption}
\usepackage{makeidx}
\usepackage{multicol}
\usepackage[all,import]{xy}
\usepackage{natbib}
\usepackage{hyperref}

\definecolor{mygray}{gray}{0.6}
\usepackage{listings}

\lstset{language=R,
    basicstyle=\small\ttfamily,
showstringspaces=false,
    commentstyle=\color{mygray},
}

\lstnewenvironment{rc}[1][]{\lstset{language=R}}{}
\newcommand{\ri}[1]{\lstinline{#1}}

\lstset{language=R}

\renewcommand{\theparagraph}{\arabic{chapter}.\arabic{paragraph}}
\setcounter{secnumdepth}{5}
\newenvironment{myproof}[2] {\noindent {\bfseries Proof of {#1} {#2}:}}{\hfill$\square$}
\newenvironment{mysolution}[1]{\vspace{0.1in} \noindent {\bfseries Solution to Problem {#1}:}}{\hfill$\square$}

\newtheorem{lemma}{Lemma}
\newtheorem{proposition}{Proposition}
\newtheorem{assumption}{Assumption}
\newtheorem{remark}{Remark}
\newtheorem{condition}{Condition}
\newtheorem{corollary}{Corollary}
\newtheorem{theorem}{Theorem}
\newtheorem{example}{Example}
\newtheorem{definition}{Definition}

\newcommand{\pkg}[1]{{\fontseries{b}\selectfont #1}}
\def\letas{\mathrel{\mathop{=}\limits^{\triangle}}}
\def\ind{\begin{picture}(9,8)
         \put(0,0){\line(1,0){9}}
         \put(3,0){\line(0,1){8}}
         \put(6,0){\line(0,1){8}}
         \end{picture}
        }
\def\nind{\begin{picture}(9,8)
         \put(0,0){\line(1,0){9}}
         \put(3,0){\line(0,1){8}}
         \put(6,0){\line(0,1){8}}
         \put(1,0){{\it /}}
         \end{picture}
    }

\def\AVar{\text{AsyVar}}
\def\var{\textup{var}}
\def\cov{\textup{cov}}
\def\sumn{\sum_{i=1}^n}
\def\summ{\sum_{j=1}^m}
\def\convergeas{\stackrel{a.s.}{\longrightarrow}}
\def\converged{\stackrel{\textup{d}}{\longrightarrow}}
\def\iidsim{\stackrel{\textup{IID}}{\sim}}
\def\indsim{\stackrel{ind}{\sim}}
\def\pr{\textup{pr}}
\def\RD{\textsc{rd}}
\def\RR{\textsc{rr}}
\def\OR{\textsc{or}}
\def\true{\textup{true}}
\def\obs{\textup{obs}}
\def\G{\text{Gamma}}
\def\d{\textnormal{d}}
\def\N01{\textsc{N}(0,1)}
\def\HF{H_{0\textsc{f}}}
\def\HN{H_{0\textsc{n}}}
\newcommand*{\tran }{^{\mkern-1.5mu\mathsf{T}}}
\def\asim{\stackrel{\cdot}{\sim}}
\def\diff{\textup{d}}
\def\at{\textup{a}}
\def\nt{\textup{n}}
\def\cp{\textup{c}}
\def\df{\textup{d}}
\def\logit{\textup{logit}}
\def\expit{\textup{expit}}
\def\L{\ell}
\def\U{u}
\def\CDE{\textsc{cde}}
\def\NDE{\textsc{nde}}
\def\NIE{\textsc{nie}}
\def\se{\textup{se}}
\def\frt{\textsc{frt}}
\def\CACE{\textsc{cace}}
\def\LATE{\textsc{late}}
\def\TSLS{\textsc{tsls}}
\def\ILS{\textsc{ils}}

\frenchspacing
\tolerance=5000

\makeindex

\title{Problem Set 5 Preparation}
\author{Name: \underline{\hspace{3cm}}}
\date{Due: May 27, 4:30 pm}

\begin{document}

\frontmatter

\maketitle
\tableofcontents

\mainmatter

\makeatletter
\@addtoreset{assumption}{chapter}
\@addtoreset{example}{chapter}
\@addtoreset{lemma}{chapter}
\@addtoreset{remark}{chapter}
\@addtoreset{theorem}{chapter}
\@addtoreset{corollary}{chapter}
\@addtoreset{proposition}{chapter}
\@addtoreset{definition}{chapter}
\makeatother

\renewcommand {\theexample} {\thechapter.\arabic{example}}
\renewcommand {\theassumption} {\thechapter.\arabic{assumption}}
\renewcommand {\thelemma} {\thechapter.\arabic{lemma}}
\renewcommand {\thetheorem} {\thechapter.\arabic{theorem}}
\renewcommand {\theremark} {\thechapter.\arabic{remark}}
\renewcommand {\thecondition} {\thechapter.\arabic{condition}}
\renewcommand {\thecorollary} {\thechapter.\arabic{corollary}}
\renewcommand {\theproposition} {\thechapter.\arabic{proposition}}
\renewcommand {\thedefinition} {\thechapter.\arabic{definition}}

\chapter{Problem Set 5 Preparation}

\section{Notation and Basic Results}

This document collects the problem statements, notation, and relevant results
needed for Problem Set 5. It follows the style of Ding's textbook: state the
setup, recall the assumptions and results, and leave the solution below each
problem. The solution and proof parts are intentionally left blank.

Consider units indexed by \(i=1,\ldots,n\). Let \(Z_i\) be the treatment assigned
or instrumental variable, \(D_i\) be the treatment received, and \(Y_i\) be the
outcome. For a binary \(D\), define the latent compliance type
\[
U_i=\left\{
\begin{array}{ll}
\at, & D_i(1)=1,\ D_i(0)=1,\\
\cp, & D_i(1)=1,\ D_i(0)=0,\\
\df, & D_i(1)=0,\ D_i(0)=1,\\
\nt, & D_i(1)=0,\ D_i(0)=0.
\end{array}
\right.
\]
Here \(\at,\cp,\df,\nt\) denote always takers, compliers, defiers, and never
takers, respectively.

\begin{assumption}[randomization]\label{assume::random-hw}
\[
Z\ind \{D(1),D(0),Y(1),Y(0)\}.
\]
\end{assumption}

\begin{assumption}[monotonicity]\label{assume::Mono-hw}
\[
\pr(U=\df)=0,
\]
or equivalently \(D_i(1)\geq D_i(0)\) for all \(i\).
\end{assumption}

\begin{assumption}[exclusion restriction]\label{assume::ER-hw}
\[
Y_i(1)=Y_i(0)
\]
for always takers and never takers.
\end{assumption}

\begin{definition}[CACE or LATE]\label{def::CACE-hw}
Define
\[
\tau_\cp
=
E\{Y(1)-Y(0)\mid U=\cp\}
\]
as the complier average causal effect or local average treatment effect.
\end{definition}

\begin{theorem}[identification of CACE]\label{thm::CACE-identification-hw}
Under Assumptions \ref{assume::random-hw}--\ref{assume::ER-hw},
\[
\tau_\cp
=
\frac{E(Y\mid Z=1)-E(Y\mid Z=0)}
{E(D\mid Z=1)-E(D\mid Z=0)}.
\]
The corresponding sample ratio is the Wald estimator or IV estimator
\[
\hat\tau_\cp=\frac{\hat\tau_Y}{\hat\tau_D}.
\]
\end{theorem}

\begin{remark}
When the potential outcome is indexed by both assignment and treatment received,
the exclusion restriction is written as \(Y_i(z,d)=Y_i(d)\). In that notation,
\[
\tau_\cp=E\{Y(d=1)-Y(d=0)\mid U=\cp\}.
\]
\end{remark}

\section{Homework problems}

\paragraph{Variance of the Wald estimator}\label{problem::infinite-var-iv-hw}

Show that
\[
\var(\hat{\tau}_\cp)=\infty.
\]

Source: Ding, Chapter 21, Problem 21.1, corresponding to
\ri{problem::infinite-var-iv}.

The following facts are useful. The Wald estimator has the form
\[
\hat{\tau}_\cp
=
\frac{\hat{\tau}_Y}{\hat{\tau}_D},
\quad
\hat{\tau}_Y=\bar Y_1-\bar Y_0,
\quad
\hat{\tau}_D=\bar D_1-\bar D_0.
\]
Even if \(\tau_D>0\), the denominator \(\hat\tau_D\) may be zero with positive
probability in finite samples. This is the key reason that the finite-sample
variance of the ratio estimator can be infinite.

\begin{remark}
This problem is connected to the weak IV discussion in Chapter 21. The Wald
estimator was introduced by \citet{wald1940fitting}; the CACE/LATE framework is
based on \citet{imbens1994identification} and \citet{angrist1996identification}.
\end{remark}

\begin{mysolution}{21.1}

\end{mysolution}

\paragraph{Binary IV and ordinal treatment received}\label{para::iv-ordinal-treatment-hw}

\citet{angrist1995two} discussed a more general setting with a binary IV \(Z\),
an ordinal treatment received \(D\in\{0,1,\ldots,J\}\), and an outcome \(Y\).
The ordinal treatment received has potential outcomes \(D(1)\) and \(D(0)\),
and the outcome has potential outcomes \(Y(z,d)\). The relevant assumptions are
\[
Z\ind \{D(z),Y(z,d):z=0,1;\ d=0,1,\ldots,J\},
\]
\[
D(1)\geq D(0),
\quad
Y(z,d)=Y(d).
\]
Prove that
\[
\frac{E(Y\mid Z=1)-E(Y\mid Z=0)}
{E(D\mid Z=1)-E(D\mid Z=0)}
=
\sum_{j=1}^J w_j
E\{Y(j)-Y(j-1)\mid D(1)\geq j>D(0)\},
\]
where
\[
w_j=
\frac{\pr\{D(1)\geq j>D(0)\}}
{\sum_{j'=1}^J\pr\{D(1)\geq j'>D(0)\}}.
\]

Source: Ding, Chapter 21, Problem 21.6, corresponding to
\ri{para::iv-ordinal-treatment}.

When \(J=1\), the result reduces to Theorem
\ref{thm::CACE-identification-hw}. A useful starting point is
\[
D(z)=\sum_{j=0}^J j1\{D(z)=j\},
\quad
Y(D(z))=\sum_{j=0}^J Y(j)1\{D(z)=j\}.
\]
Ding suggests using Abel's lemma, also called summation by parts:
\[
\sum_{j=0}^J f_j(g_{j+1}-g_j)
=
f_Jg_{J+1}-f_0g_0-\sum_{j=1}^J g_j(f_j-f_{j-1}).
\]

\begin{remark}
The target estimand is a weighted average of adjacent-level causal effects for
latent groups defined by \(D(1)\geq j>D(0)\).
\end{remark}

\begin{mysolution}{21.6}

\end{mysolution}

\paragraph{Data analysis: a flu shot encouragement design}\label{problem::flushot-hw}

The dataset \ri{fludata.txt} is from a randomized encouragement design of
\citet{mcdonald1992effects}, which was also re-analyzed by
\citet{hirano2000assessing}. It contains the following variables:
\[
\begin{array}{ll}
\hline
\texttt{assign} & \text{binary encouragement to receive the flu shot}\\
\texttt{receive} & \text{binary indicator for receiving the flu shot}\\
\texttt{outcome} & \text{binary outcome for flu-related hospitalization}\\
\texttt{age} & \text{age of the patient}\\
\texttt{sex} & \text{sex of the patient}\\
\texttt{race} & \text{race of the patient}\\
\texttt{copd} & \text{chronic obstructive pulmonary disease}\\
\texttt{dm} & \text{diabetes}\\
\texttt{heartd} & \text{heart disease}\\
\texttt{renal} & \text{renal disease}\\
\texttt{liverd} & \text{liver disease}\\
\hline
\end{array}
\]
Analyze the data with and without adjusting for the covariates.

Source: Ding, Chapter 21, Problem 21.7, corresponding to
\ri{problem::flushot}. The local data file is
\ri{/Users/yitwah/homework/hw-causal/fludata.txt}.

Without covariates, use the Wald estimator
\[
\hat\tau_\cp
=
\frac{\bar Y_{Z=1}-\bar Y_{Z=0}}
{\bar D_{Z=1}-\bar D_{Z=0}}.
\]
With covariates in a completely randomized experiment, use Lin-style covariate
adjustment for both \(Y\) and \(D\), then take the ratio
\[
\hat\tau_{\cp,L}
=
\frac{\hat\tau_{Y,L}}{\hat\tau_{D,L}}.
\]

\begin{remark}
Report the first stage \(\hat\tau_D\), ITT effect \(\hat\tau_Y\), CACE estimate,
standard error or confidence interval, and the assumptions used for the causal
interpretation.
\end{remark}

\begin{mysolution}{21.7}

\end{mysolution}

\paragraph{Risk ratio for compliers}\label{hw::rr-compliers-hw}

With a binary outcome, define the risk ratio for compliers as
\[
\RR_\cp
=
\frac{\pr\{Y(1)=1\mid U=\cp\}}
{\pr\{Y(0)=1\mid U=\cp\}}.
\]
Show that under Assumptions \ref{assume::random-hw}--\ref{assume::ER-hw},
\[
\RR_\cp
=
\frac{E(DY\mid Z=1)-E(DY\mid Z=0)}
{E\{(D-1)Y\mid Z=1\}-E\{(D-1)Y\mid Z=0\}}.
\]

Source: Ding, Chapter 22, Problem 22.2, corresponding to
\ri{hw::rr-compliers}.

We first recall the mixture identification formulas:
\[
\pi_\nt=\pr(D=0\mid Z=1),
\quad
\pi_\at=\pr(D=1\mid Z=0),
\quad
\pi_\cp=E(D\mid Z=1)-E(D\mid Z=0).
\]
The complier potential outcome means are
\[
\mu_{1\cp}
=
\pi_\cp^{-1}\{E(DY\mid Z=1)-E(DY\mid Z=0)\},
\]
and
\[
\mu_{0\cp}
=
\pi_\cp^{-1}
\left[E\{(1-D)Y\mid Z=0\}-E\{(1-D)Y\mid Z=1\}\right].
\]

\begin{remark}
Since \((D-1)Y=-(1-D)Y\), the denominator in the target expression equals the
identified numerator of \(\mu_{0\cp}\).
\end{remark}

\begin{mysolution}{22.2}

\end{mysolution}

\paragraph{Bounds on the average causal effect on the whole population}\label{para::bounds-ACE-hw}

Use the potential outcome \(Y(d)\) and define the average causal effect of the
treatment received on the whole population as
\[
\delta=E\{Y(d=1)-Y(d=0)\}.
\]
Let
\[
m_{du}=E\{Y(d)\mid U=u\},\quad u=\at,\nt,\cp.
\]
Then
\[
\delta=\sum_{u=\at,\nt,\cp}\pi_u(m_{1u}-m_{0u}).
\]
The data identify
\(\pi_\at,\pi_\nt,\pi_\cp,m_{1\at},m_{0\nt},m_{1\cp},m_{0\cp}\), but not
\(m_{0\at}\) and \(m_{1\nt}\). With \(Y\in[\underline y,\overline y]\), prove
\[
\underline\delta\leq \delta\leq \overline\delta,
\]
where
\[
\underline{\delta}
=
\delta'
-\overline y\,\pr(D=1\mid Z=0)
+\underline y\,\pr(D=0\mid Z=1),
\]
\[
\overline{\delta}
=
\delta'
-\underline y\,\pr(D=1\mid Z=0)
+\overline y\,\pr(D=0\mid Z=1),
\]
and
\[
\delta'=E(DY\mid Z=1)-E(Y-DY\mid Z=0).
\]

Source: Ding, Chapter 22, Problem 22.7, corresponding to
\ri{para::bounds-ACE}.

The useful identified proportions are
\[
\pi_\at=\pr(D=1\mid Z=0),
\quad
\pi_\nt=\pr(D=0\mid Z=1),
\quad
\pi_\cp=E(D\mid Z=1)-E(D\mid Z=0).
\]
The unidentified terms are bounded only by the outcome support:
\[
m_{0\at}\in[\underline y,\overline y],
\quad
m_{1\nt}\in[\underline y,\overline y].
\]

\begin{remark}
For a binary outcome, Ding gives the simplified bounds
\[
\underline{\delta}
=
E(DY\mid Z=1)-E(D+Y-DY\mid Z=0),
\]
and
\[
\overline{\delta}
=
E(DY+1-D\mid Z=1)-E(Y-DY\mid Z=0).
\]
\end{remark}

\begin{mysolution}{22.7}

\end{mysolution}

\paragraph{More algebra for TSLS}\label{prob::algebra-tsls-hw}

Show that the \(\hat\Gamma\) in \(\hat D_i=\hat\Gamma\tran Z_i\) equals
\[
\hat{\Gamma}
=
\left(\sumn Z_iZ_i\tran\right)^{-1}\sumn Z_iD_i\tran.
\]
Then show that \(\hat\beta_\TSLS\) reduces to \(\hat\beta_{\textsc{iv}}\) if
\(Z\) and \(D\) have the same dimension and
\[
n^{-1}\sumn Z_iZ_i\tran,
\quad
n^{-1}\sumn Z_iD_i\tran
\]
are both invertible.

Source: Ding, Chapter 23, Problem 23.1, corresponding to
\ri{prob::algebra-tsls}.

The two-stage least squares estimator \citep{theil1953estimation,basmann1957generalized}
is defined by first regressing \(D\) on \(Z\) to get \(\hat D\), and then
regressing \(Y\) on \(\hat D\). Explicitly,
\[
\hat\beta_\TSLS
=
\left(\sumn \hat D_i\hat D_i\tran\right)^{-1}\sumn \hat D_iY_i.
\]
The first-stage OLS fit gives \(D_i=\hat D_i+\check D_i\) and the orthogonality
\[
\sumn \hat D_i\check D_i\tran=0.
\]

\begin{remark}
In the just-identified case, the relevant IV estimator is
\[
\hat\beta_{\textsc{iv}}
=
\left(\sumn Z_iD_i\tran\right)^{-1}\sumn Z_iY_i.
\]
\end{remark}

\begin{mysolution}{23.1}

\end{mysolution}

\paragraph{Equivalence between TSLS and ILS}\label{para::tsls-ils-hw}

Prove that, with a single endogenous treatment and a single IV,
\[
\hat\beta_{1,\ILS}=\hat\beta_{1,\TSLS}.
\]

Source: Ding, Chapter 23, Problem 23.2, corresponding to
\ri{para::tsls-ils}. This problem is optional in Problem Set 5.

The structural form is
\[
\left\{
\begin{array}{l}
Y_i=\beta_0+\beta_1D_i+\beta_2\tran X_i+\varepsilon_i,\\
D_i=\gamma_0+\gamma_1Z_i+\gamma_2\tran X_i+\varepsilon_{2i}.
\end{array}
\right.
\]
The reduced form is
\[
\left\{
\begin{array}{l}
Y_i=\Gamma_0+\Gamma_1Z_i+\Gamma_2\tran X_i+\varepsilon_{1i},\\
D_i=\gamma_0+\gamma_1Z_i+\gamma_2\tran X_i+\varepsilon_{2i}.
\end{array}
\right.
\]
The ILS estimator is
\[
\hat\beta_{1,\ILS}=\frac{\hat\Gamma_1}{\hat\gamma_1}.
\]

\begin{remark}
Ding suggests using the FWL theorem. The equality was noted by
\citet[][Section A.3]{imbens2014instrumental}.
\end{remark}

\begin{mysolution}{23.2}

\end{mysolution}

\paragraph{Control function in the linear instrumental variable model}\label{para::control-function-iv-hw}

Define the control function estimator \(\hat\beta_{\textsc{cf}}\) as follows.
First run OLS of \(D\) on \(Z\), and obtain residuals \(\check D_i\). Then run
OLS of \(Y\) on \(D\) and \(\check D\), and take the coefficient of \(D\). Show
that
\[
\hat\beta_{\textsc{cf}}=\hat\beta_\TSLS.
\]

Source: Ding, Chapter 23, Problem 23.3, corresponding to
\ri{para::control-function-iv}.

Use the first-stage decomposition
\[
D=\hat D+\check D.
\]
The fitted value \(\hat D\) is orthogonal to the residual \(\check D\). The
FWL theorem and OLS orthogonality are the main algebraic tools.

\begin{remark}
\citet{hausman1978specification} pointed out this result, and
\citet{wooldridge2015control} provides a broader discussion of control function
methods.
\end{remark}

\begin{mysolution}{23.3}

\end{mysolution}

\section{Additional problems from Problem Set 5}

\paragraph{Wald and two-stage least squares}\label{hw::wald-cov-form}

We derived the identification formula for CACE under randomization in class.
Show that
\[
\CACE
=
\frac{\cov(Z_i,Y_i)}{\cov(Z_i,D_i)}.
\]

The starting point is Theorem \ref{thm::CACE-identification-hw}. For a binary
\(Z\), with \(p=\pr(Z=1)\),
\[
\cov(Z,Y)=p(1-p)\{E(Y\mid Z=1)-E(Y\mid Z=0)\},
\]
and
\[
\cov(Z,D)=p(1-p)\{E(D\mid Z=1)-E(D\mid Z=0)\}.
\]
Chapter 23 gives the same formula from the scalar IV model:
\[
\beta=\frac{\cov(Z,Y)}{\cov(Z,D)}.
\]

\begin{remark}
This problem links the potential-outcome CACE formula to the linear IV/TSLS
notation.
\end{remark}

\begin{mysolution}{2}

\end{mysolution}

\paragraph{Weighting method for CACE in observational studies}\label{hw::weighting-cace-x}

Suppose randomization is replaced by conditional ignorability
\[
Z_i\ind \{Y_i(1),Y_i(0),D_i(1),D_i(0)\}\mid X_i,
\]
while monotonicity and exclusion restriction continue to hold. Define
\[
\kappa_{0i}
=
(1-D_i)
\frac{(1-Z_i)-\pr(Z_i=0\mid X_i)}
{\pr(Z_i=0\mid X_i)\pr(Z_i=1\mid X_i)},
\]
and
\[
\kappa_{1i}
=
D_i
\frac{Z_i-\pr(Z_i=1\mid X_i)}
{\pr(Z_i=0\mid X_i)\pr(Z_i=1\mid X_i)}.
\]
Show that
\[
E\{g(Y_i(0),X_i)\mid U_i=\cp\}
=
\frac{1}{\pr(U_i=\cp)}E\{\kappa_{0i}g(Y,X)\},
\]
and
\[
E\{g(Y_i(1),X_i)\mid U_i=\cp\}
=
\frac{1}{\pr(U_i=\cp)}E\{\kappa_{1i}g(Y,X)\}.
\]

Source: Problem Set 5, Problem 3. This problem is optional.

Let \(e(X)=\pr(Z=1\mid X)\). Conditional ignorability allows observed conditional
expectations within \(Z\)-groups to represent conditional potential-outcome
expectations. The monotonicity and exclusion restriction assumptions keep the
same latent compliance interpretation as in Chapter 21.

\begin{remark}
This problem is the observational-study analogue of the mixture disentangling
identities in Chapter 22.
\end{remark}

\begin{mysolution}{3}

\end{mysolution}

\paragraph{Analysis of return to schooling}\label{hw::card-schooling}

\citet{card1993using} used geographic variation in college proximity to
estimate the return to schooling. The treatment is education, the outcome is log
wage, and the IV is college proximity.

First, dichotomize the education variable and infer the CACE for units with
\(D(1)=1,D(0)=0\), where \(D\) is the dichotomized education. Discuss how to
choose the threshold and whether the results are sensitive to the choice. Conduct
an analysis with covariates. Second, without dichotomizing education, use the
standard TSLS estimator and compare the results.

Source: Problem Set 5, Problem 4. The same Card data-analysis problem appears as
a commented problem in Ding, Chapter 23.

For a threshold \(c\), define \(D_c=1\{\text{education}\geq c\}\). The
dichotomized causal quantity is
\[
\tau_\cp(c)
=
E\{Y(1)-Y(0)\mid D_c(1)=1,D_c(0)=0\}.
\]
Without covariates,
\[
\hat\tau_\cp(c)
=
\frac{\bar Y_{Z=1}-\bar Y_{Z=0}}
{\bar D_{c,Z=1}-\bar D_{c,Z=0}}.
\]
With covariates, adjust both \(Y\) and \(D_c\), then take the ratio of the
adjusted effects. Without dichotomization, the TSLS specification is
\[
\left\{
\begin{array}{l}
\text{education}_i=\gamma_0+\gamma_1Z_i+\gamma_2\tran X_i+\varepsilon_{2i},\\
\log(\text{wage}_i)=\beta_0+\beta_1\text{education}_i+\beta_2\tran X_i+\varepsilon_i.
\end{array}
\right.
\]

\begin{remark}
The current folder does not contain \ri{card.csv} or \ri{code\_bk.txt}.
Add their local paths here before doing the empirical analysis.
\end{remark}

\begin{mysolution}{4}

\end{mysolution}

\bibliographystyle{plainnat}
\bibliography{causal/causal}

\end{document}
